\documentclass[11pt,a4paper]{article}

% ==========================================================
% Configuracion general
% ==========================================================
\usepackage[utf8]{inputenc}
\usepackage[T1]{fontenc}
\usepackage{lmodern}
\usepackage{microtype}
\usepackage[spanish]{babel}
\usepackage[margin=2.2cm]{geometry}
\usepackage{graphicx}
\usepackage{booktabs}
\usepackage{array}
\usepackage{tabularx}
\usepackage{longtable}
\usepackage{amssymb}
\usepackage{tikz}
\usetikzlibrary{arrows.meta,positioning}
\usepackage{float}
\usepackage{hyperref}
\usepackage{xcolor}
\usepackage{enumitem}
\usepackage{titlesec}
\usepackage{fancyhdr}

% Paleta de color (misma familia que el Avance 1 del Proyecto Integrador)
\definecolor{myNewColorA}{RGB}{27,94,55}
\definecolor{myNewColorB}{RGB}{46,125,80}
\definecolor{myNewColorC}{RGB}{184,218,191}

\titleformat{\section}{\Large\bfseries\color{myNewColorA}}{\thesection.}{0.5em}{}
\titleformat{\subsection}{\large\bfseries\color{myNewColorB}}{\thesubsection}{0.5em}{}

\hypersetup{
    colorlinks=true,
    linkcolor=myNewColorA,
    urlcolor=myNewColorA,
    citecolor=myNewColorA,
    pdftitle={Practica 2 - Taller SQL Relacional - Grupo 04},
    pdfauthor={Malan Mullo, Marco Vinicio; Pujos Culque, Viviana Isabel; Chuncho Morocho, Carlos Guillermo}
}

\pagestyle{fancy}
\fancyhf{}
\renewcommand{\headrulewidth}{0.4pt}
\fancyhead[L]{\small Practica 2 -- Taller SQL Relacional}
\fancyhead[R]{\small M1721 -- Grupo 04}
\fancyfoot[C]{\small\thepage}

\setlist[itemize]{itemsep=1pt, topsep=2pt}
\setlist[enumerate]{itemsep=1pt, topsep=2pt}

\newcolumntype{L}[1]{>{\raggedright\arraybackslash}p{#1}}
\newcolumntype{C}[1]{>{\centering\arraybackslash}p{#1}}

\emergencystretch=3em

\begin{document}

% ==========================================================
% PORTADA
% ==========================================================
\begin{titlepage}
\centering
\vspace*{1cm}

{\large ESCUELA SUPERIOR POLIT\'ECNICA DE CHIMBORAZO (ESPOCH)\par}
\vspace{0.3cm}
{\large Maestr\'ia en Estad\'istica con menci\'on en Ciencia de Datos e Inteligencia Artificial\par}

\vspace{2.5cm}
{\Huge\bfseries\color{myNewColorA} Pr\'actica 2 -- Taller SQL\par}
\vspace{0.5cm}
{\large Modelado, normalizaci\'on y consultas\par}

\vspace{1.5cm}
\begin{minipage}{0.85\textwidth}
\centering
{\Large\bfseries Caso aplicado\par}
\vspace{0.3cm}
{\large Detecciones de incendio, precipitaci\'on y cobertura vegetal por celda espacial\\
y fecha, en el cant\'on Loja\par}
\end{minipage}

\vspace{2cm}
{\large\bfseries Integrantes\par}
\vspace{0.3cm}
{\large
MALAN MULLO MARCO VINICIO \\
PUJOS CULQUE VIVIANA ISABEL \\
CARLOS GUILLERMO CHUNCHO M. \par}

\vspace{2cm}
{\large\bfseries Grupo 04}\par

\vspace{1.5cm}
{\large Pr\'actica 2 -- M1721 \par}
{\large 2026 \par}

\vfill
\end{titlepage}

% ==========================================================
% 1. CASO APLICADO
% ==========================================================
\section{Caso aplicado}

\textbf{Contexto.} Las detecciones satelitales de incendio, los registros de precipitaci\'on y la
cobertura vegetal de una zona suelen analizarse por separado, sin relacionarlos entre s\'i por
celda espacial y fecha. Este taller usa el mismo dominio tem\'atico que el Proyecto Integrador
FireForest (incendios forestales, cant\'on Loja), pero como una pr\'actica \textbf{independiente},
con su propio esquema, sus propios datos de ejemplo y su propia base de datos PostgreSQL
(\texttt{taller\_sql\_fireforest}).

\textbf{Problema.} Sin una estructura relacional integrada, preguntas simples como "¿qu\'e
celdas con cierto tipo de cobertura vegetal concentran m\'as detecciones?" o "¿hay meses con
varias detecciones y poca lluvia?" requieren cruzar manualmente fuentes distintas.

\textbf{Prop\'osito de la base de datos.} Diseñar e implementar una base de datos relacional en
PostgreSQL que integre, sin redundancia, detecciones de incendio, precipitaci\'on diaria y
cobertura vegetal por celda y fecha, y que pueda consultarse mediante SQL para responder
preguntas operativas y anal\'iticas descriptivas.

\textbf{Usuarios potenciales.} Estudiantes y docentes del curso, como ejercicio de modelado y
SQL; y, como referencia conceptual, personal t\'ecnico ambiental que quisiera una primera
aproximaci\'on relacional (sin que exista convenio real con ninguna instituci\'on).

\textbf{Preguntas operativas.}
\begin{enumerate}[itemsep=0pt]
    \item ¿Qu\'e detecciones de incendio existen para una celda y un periodo espec\'ificos?
    \item ¿Cu\'al fue la precipitaci\'on registrada para una celda y una fecha determinadas?
    \item ¿Qu\'e celdas no tienen registro de precipitaci\'on para una fecha con incendio
          observado?
\end{enumerate}

\textbf{Preguntas anal\'iticas.}
\begin{enumerate}[itemsep=0pt]
    \item ¿Cu\'antas detecciones y qu\'e FRP promedio/m\'aximo presenta cada celda?
    \item ¿Qu\'e celdas y meses acumulan precipitaci\'on mensual baja?
    \item ¿Qu\'e combinaciones celda-mes presentan a la vez detecciones y precipitaci\'on baja?
    \item ¿Qu\'e tipo de cobertura vegetal concentra m\'as detecciones? (usa la relaci\'on N:M)
\end{enumerate}

\textit{Nota: todos los datos usados en este taller son de ejemplo, controlados, generados
espec\'ificamente para esta pr\'actica. No son datos reales de VIIRS ni de CHIRPS.}

% ==========================================================
% 2. MODELO RELACIONAL
% ==========================================================
\section{Modelo relacional}

\textbf{Entidades principales:} celda espacial (\texttt{dim\_celda}), fecha
(\texttt{dim\_fecha}), cobertura vegetal (\texttt{dim\_cobertura\_vegetal}), detecci\'on de
incendio (\texttt{fact\_incendio}) y registro de precipitaci\'on (\texttt{fact\_clima}).

\textbf{Relaciones 1:N:} \texttt{dim\_celda} con \texttt{fact\_incendio} y con
\texttt{fact\_clima}; \texttt{dim\_fecha} con \texttt{fact\_incendio} y con \texttt{fact\_clima}.

\textbf{Relaci\'on N:M:} una celda puede tener m\'as de un tipo de cobertura vegetal registrado
en un mismo a\~no (mosaico de uso de suelo), y un tipo de cobertura aplica a muchas celdas. Esta
relaci\'on muchos-a-muchos entre \texttt{dim\_celda} y \texttt{dim\_cobertura\_vegetal} se
resuelve con la tabla intermedia \texttt{celda\_cobertura(celda\_id, cobertura\_id, anio)}.

\begin{center}
\scriptsize
\begin{tabular}{@{}L{3.2cm}L{3.8cm}L{1.8cm}L{2.6cm}L{3.4cm}@{}}
\toprule
\textbf{Tabla} & \textbf{Prop\'osito} & \textbf{PK} & \textbf{FK} & \textbf{Cardinalidad} \\
\midrule
\texttt{dim\_celda} & Celda espacial de referencia (500 m $\times$ 500 m) & \texttt{celda\_id} & --- & 1 celda -- N incendio/clima/cobertura \\
\texttt{dim\_fecha} & Calendario diario del periodo & \texttt{fecha\_id} & --- & 1 fecha -- N incendio/clima \\
\texttt{dim\_cobertura\_vegetal} & Cat\'alogo de tipos de cobertura & \texttt{cobertura\_id} & --- & 1 cobertura -- N celda\_cobertura \\
\texttt{celda\_cobertura} & Relaci\'on N:M celda$\leftrightarrow$cobertura por a\~no & \texttt{(celda\_id, cobertura\_id, anio)} & celda\_id, cobertura\_id & N:M \\
\texttt{fact\_incendio} & Detecci\'on diaria de incendio por celda & \texttt{incendio\_id} & celda\_id, fecha\_id & N -- 1 celda; N -- 1 fecha \\
\texttt{fact\_clima} & Precipitaci\'on diaria por celda & \texttt{clima\_id} & celda\_id, fecha\_id & N -- 1 celda; N -- 1 fecha \\
\bottomrule
\end{tabular}
\end{center}

\begin{figure}[H]
\centering
\begin{tikzpicture}[
  box/.style={draw=myNewColorA, rounded corners, fill=myNewColorC!30,
              align=left, text width=3.5cm, font=\scriptsize, inner sep=3pt},
  arr/.style={-{Stealth[length=1.8mm]}, thick, myNewColorA},
  lbl/.style={font=\scriptsize, fill=white, inner sep=1pt}
]
\node[box] (celda) at (0,4) {\textbf{dim\_celda}\\ PK: celda\_id};
\node[box] (cobertura) at (7,4) {\textbf{dim\_cobertura\_vegetal}\\ PK: cobertura\_id};
\node[box] (cc) at (3.5,2) {\textbf{celda\_cobertura}\\ PK: (celda\_id,\\ cobertura\_id, anio)\\ FK: ambas};
\node[box] (fecha) at (0,0) {\textbf{dim\_fecha}\\ PK: fecha\_id};
\node[box] (incendio) at (-2.5,-2.2) {\textbf{fact\_incendio}\\ PK: incendio\_id\\ FK: celda\_id, fecha\_id};
\node[box] (clima) at (2.5,-2.2) {\textbf{fact\_clima}\\ PK: clima\_id\\ FK: celda\_id, fecha\_id};

\draw[arr] (celda) -- (cc) node[lbl, midway] {1:N};
\draw[arr] (cobertura) -- (cc) node[lbl, midway] {1:N};
\draw[arr] (celda) -- (incendio) node[lbl, pos=0.35] {1:N};
\draw[arr] (celda) -- (clima) node[lbl, pos=0.35] {1:N};
\draw[arr] (fecha) -- (incendio) node[lbl, midway] {1:N};
\draw[arr] (fecha) -- (clima) node[lbl, midway] {1:N};
\end{tikzpicture}
\caption{Modelo relacional final. La relaci\'on N:M entre \texttt{dim\_celda} y
\texttt{dim\_cobertura\_vegetal} queda resuelta por \texttt{celda\_cobertura}.}
\label{fig:modelo}
\end{figure}

% ==========================================================
% 3. NORMALIZACION
% ==========================================================
\section{Normalizaci\'on}

\textbf{Modelo inicial (desnormalizado).} Antes de dividir en tablas, la informaci\'on cabr\'ia
en un registro plano por celda-d\'ia, con columnas repetidas para cobertura vegetal
(\texttt{cobertura\_1/pct\_1}, \texttt{cobertura\_2/pct\_2}, ...) porque una celda puede tener
m\'as de un tipo de cobertura.

\begin{figure}[H]
\centering
\begin{tikzpicture}[
  box/.style={draw=myNewColorA, rounded corners, fill=myNewColorC!20,
              align=left, text width=8.2cm, font=\scriptsize, inner sep=5pt},
  rep/.style={draw=myNewColorB, dashed, rounded corners, inner sep=3pt}
]
\node[box] (reg) {
  \textbf{REGISTRO\_DIARIO} (sin PK definida; una fila por celda-d\'ia)\\[2pt]
  celda\_id, longitud, latitud, area\_km2\\
  fecha, anio, mes\\[4pt]
  \tikz[baseline]{\node[rep, inner sep=3pt] {cobertura\_1, pct\_1, cobertura\_2, pct\_2, \ldots};}\\[2pt]
  {\scriptsize\color{myNewColorB} (grupo repetitivo, de longitud variable)}\\[4pt]
  numero\_detecciones, frp\_suma\_mw, frp\_maxima\_mw, incendio\_observado\\
  precipitacion\_diaria\_mm
};
\end{tikzpicture}
\caption{Modelo inicial (desnormalizado): una \'unica entidad plana, con un grupo repetitivo de
cobertura vegetal que impide una PK simple y viola 1FN.}
\label{fig:modelo-inicial}
\end{figure}

\textbf{1FN.} Se eliminan los grupos repetitivos de cobertura extray\'endolos a
\texttt{celda\_cobertura} (una fila por combinaci\'on celda-cobertura-a\~no). Todos los atributos
del esquema final son at\'omicos (sin listas ni valores compuestos) y cada tabla tiene PK
expl\'icita.

\textbf{2FN.} Las tablas con PK simple no pueden tener dependencia parcial. La \'unica tabla con
PK compuesta es \texttt{celda\_cobertura(celda\_id, cobertura\_id, anio)}, donde
$(\texttt{celda\_id}, \texttt{cobertura\_id}, \texttt{anio}) \rightarrow
\texttt{porcentaje\_cobertura}$: el porcentaje depende de los tres atributos a la vez (el mismo
tipo de cobertura puede tener otro porcentaje en otra celda o en otro a\~no), sin dependencia
parcial.

\textbf{3FN.} Por tabla: \texttt{dim\_celda}: \texttt{celda\_id} $\rightarrow$
\texttt{nombre\_referencia, longitud, latitud, area\_km2}; \texttt{fact\_incendio}:
\texttt{incendio\_id} $\rightarrow$ (celda\_id, fecha\_id, fuente, numero\_detecciones,
frp\_suma\_mw, frp\_maxima\_mw, fire\_mask\_maximo, incendio\_observado); \texttt{fact\_clima}:
\texttt{clima\_id} $\rightarrow$ (celda\_id, fecha\_id, precipitacion\_diaria\_mm). En ninguna de
estas tablas un atributo no clave determina a otro atributo no clave, por lo que cumplen 3FN.

\textbf{Excepci\'on deliberada en \texttt{dim\_fecha}.} En esta tabla,
$\texttt{fecha\_id} \rightarrow \texttt{fecha} \rightarrow \texttt{anio}, \texttt{mes}$: una
fecha calendario determina su a\~no y mes, pero \texttt{anio}/\texttt{mes} no dependen
directamente de la PK sin pasar por \texttt{fecha}. Es, en sentido estricto, una dependencia
transitiva. Se conserva de forma deliberada porque es el patr\'on est\'andar de una dimensi\'on
de tiempo (evita recalcular \texttt{EXTRACT()} en cada consulta de agregaci\'on mensual), y su
riesgo de inconsistencia queda controlado por el \texttt{CHECK chk\_fecha\_id\_consistente} y por
la restricci\'on \texttt{UNIQUE} sobre \texttt{fecha}.

% ==========================================================
% 4. IMPLEMENTACION SQL
% ==========================================================
\section{Implementaci\'on SQL}

El esquema se implement\'o en PostgreSQL con las 6 tablas descritas (\texttt{dim\_celda},
\texttt{dim\_fecha}, \texttt{dim\_cobertura\_vegetal}, \texttt{celda\_cobertura},
\texttt{fact\_incendio}, \texttt{fact\_clima}), todas con clave primaria expl\'icita y claves
for\'aneas donde corresponde. El script completo est\'a en
\texttt{Grupo04\_Practica2\_SQL.sql}; a continuaci\'on, un fragmento representativo:

\begin{scriptsize}
\begin{verbatim}
CREATE TABLE fact_incendio (
    incendio_id         VARCHAR(20) PRIMARY KEY,
    celda_id            VARCHAR(20) NOT NULL REFERENCES dim_celda(celda_id),
    fecha_id            INTEGER NOT NULL REFERENCES dim_fecha(fecha_id),
    numero_detecciones  INTEGER NOT NULL DEFAULT 0 CHECK (numero_detecciones >= 0),
    incendio_observado  BOOLEAN NOT NULL DEFAULT FALSE,
    CONSTRAINT uq_incendio_celda_fecha UNIQUE (celda_id, fecha_id)
);
\end{verbatim}
\end{scriptsize}

En conjunto, el esquema usa: \texttt{PRIMARY KEY} (simple y compuesta), \texttt{FOREIGN KEY}
(\texttt{REFERENCES}), \texttt{UNIQUE} (coordenadas de celda, fecha, nombre de cobertura,
\texttt{(celda\_id, fecha\_id)}), \texttt{CHECK} (rango de mes, valores no negativos, porcentaje
de cobertura, coherencia \texttt{fecha\_id}$\leftrightarrow$\texttt{fecha}, coherencia
observado/detecciones), \texttt{NOT NULL} en los atributos de negocio, y \texttt{DEFAULT} (\'area
por defecto, fuente \texttt{'VIIRS'}, \texttt{incendio\_observado = FALSE}).

Los datos de ejemplo son controlados: \textbf{82 registros} distribuidos en las 6 tablas (6
celdas, 12 fechas, 5 tipos de cobertura, 11 filas de \texttt{celda\_cobertura}, 18 registros de
incendio y 30 de clima), con variedad de FRP, precipitaci\'on y d\'ias con y sin detecci\'on.

% ==========================================================
% 5. CONSULTAS SQL
% ==========================================================
\section{Consultas SQL}

Las 8 consultas obligatorias y una adicional, ejecutadas contra \texttt{taller\_sql\_fireforest}.
Resultados y consultas completas verificables en \texttt{05\_resultados/resultados\_consultas.md}
y en \texttt{Grupo04\_Practica2\_SQL.sql}.

\subsection*{Consulta 1 -- SELECT + WHERE}
\textbf{Pregunta:} ¿Qu\'e detecciones de incendio existen para LJ\_001 entre julio y agosto de
2023?
\begin{scriptsize}
\begin{verbatim}
SELECT fi.incendio_id, fi.celda_id, df.fecha, fi.numero_detecciones, fi.frp_maxima_mw
FROM fact_incendio fi JOIN dim_fecha df ON df.fecha_id = fi.fecha_id
WHERE fi.celda_id = 'LJ_001' AND df.fecha BETWEEN '2023-07-01' AND '2023-08-31';
\end{verbatim}
\end{scriptsize}
\textbf{Resultado (2 filas):} INC\_0001, 2023-07-05, 2 detecciones, FRP m\'ax. 15.2 MW;
INC\_0002, 2023-08-15, 1 detecci\'on, FRP m\'ax. 10.1 MW. \\
\textbf{Interpretaci\'on:} LJ\_001 registra dos fechas con incendio observado en el rango
consultado; no hay otras fechas con registro para esta celda en ese periodo.

\subsection*{Consulta 2 -- SELECT + WHERE}
\textbf{Pregunta:} ¿Cu\'al fue la precipitaci\'on diaria en LJ\_003 durante septiembre de 2023?
\begin{scriptsize}
\begin{verbatim}
SELECT fc.celda_id, df.fecha, fc.precipitacion_diaria_mm
FROM fact_clima fc JOIN dim_fecha df ON df.fecha_id = fc.fecha_id
WHERE fc.celda_id = 'LJ_003' AND df.anio = 2023 AND df.mes = 9;
\end{verbatim}
\end{scriptsize}
\textbf{Resultado (1 fila):} 2023-09-20, 6.40 mm. \\
\textbf{Interpretaci\'on:} para LJ\_003 solo existe un registro de precipitaci\'on cargado en
septiembre de 2023.

\subsection*{Consulta 3 -- ORDER BY + LIMIT}
\textbf{Pregunta:} ¿Cu\'ales son las 5 detecciones con mayor FRP m\'axima del periodo?
\begin{scriptsize}
\begin{verbatim}
SELECT fi.incendio_id, fi.celda_id, df.fecha, fi.frp_maxima_mw, fi.numero_detecciones
FROM fact_incendio fi JOIN dim_fecha df ON df.fecha_id = fi.fecha_id
WHERE fi.incendio_observado = TRUE
ORDER BY fi.frp_maxima_mw DESC LIMIT 5;
\end{verbatim}
\end{scriptsize}
\textbf{Resultado (5 filas):} LJ\_005 (25.6 MW), LJ\_003 (22.1 MW), LJ\_002 (20.3 MW), LJ\_001
(15.2 MW), LJ\_006 (14.5 MW), cada una en una fecha distinta. \\
\textbf{Interpretaci\'on:} las cinco detecciones de mayor intensidad provienen de cinco celdas
distintas; ninguna concentra m\'as de un evento en este top 5.

\subsection*{Consulta 4 -- LEFT JOIN}
\textbf{Pregunta:} ¿Qu\'e celdas tienen incendio observado sin registro de precipitaci\'on en la
misma fecha?
\begin{scriptsize}
\begin{verbatim}
SELECT fi.celda_id, df.fecha, fi.numero_detecciones, fc.precipitacion_diaria_mm
FROM fact_incendio fi JOIN dim_fecha df ON df.fecha_id = fi.fecha_id
LEFT JOIN fact_clima fc ON fc.celda_id = fi.celda_id AND fc.fecha_id = fi.fecha_id
WHERE fi.incendio_observado = TRUE AND fc.clima_id IS NULL;
\end{verbatim}
\end{scriptsize}
\textbf{Resultado (1 fila):} LJ\_001, 2023-07-05, 2 detecciones, precipitaci\'on sin registro. \\
\textbf{Interpretaci\'on:} existe un \'unico caso sin correspondencia en todo el periodo; en el
resto de fechas con incendio observado s\'i hay un registro de precipitaci\'on asociado.

\subsection*{Consulta 5 -- INNER JOIN (3 tablas)}
\textbf{Pregunta:} listado legible de detecciones observadas, con nombre de celda y fecha
calendario.
\begin{scriptsize}
\begin{verbatim}
SELECT dc.nombre_referencia, df.fecha, df.mes, fi.numero_detecciones, fi.frp_maxima_mw
FROM fact_incendio fi
JOIN dim_celda dc ON dc.celda_id = fi.celda_id
JOIN dim_fecha df ON df.fecha_id = fi.fecha_id
WHERE fi.incendio_observado = TRUE ORDER BY df.fecha;
\end{verbatim}
\end{scriptsize}
\textbf{Resultado representativo (10 filas en total; primeras 3):} Sector La Toma (2023-07-05, 2
det., 15.2 MW); Sector Zamora Huayco (2023-07-15, 2 det., 14.5 MW); Sector El Pedestal
(2023-07-25, 4 det., 22.1 MW). Resultado completo en \texttt{05\_resultados/resultados\_consultas.md}. \\
\textbf{Interpretaci\'on:} las 10 fechas con incendio observado se distribuyen en los tres meses
del periodo, sin que uno solo concentre la mayor\'ia de los casos.

\subsection*{Consulta 6 -- GROUP BY + COUNT/SUM/AVG/MAX}
\textbf{Pregunta:} ¿cu\'antas detecciones totales y qu\'e FRP promedio/m\'aximo presenta cada
celda?
\begin{scriptsize}
\begin{verbatim}
SELECT dc.celda_id, COUNT(fi.incendio_id) dias, SUM(fi.numero_detecciones) detecciones,
       ROUND(AVG(fi.frp_maxima_mw),2) frp_prom, MAX(fi.frp_maxima_mw) frp_max
FROM dim_celda dc JOIN fact_incendio fi ON fi.celda_id = dc.celda_id
GROUP BY dc.celda_id ORDER BY detecciones DESC;
\end{verbatim}
\end{scriptsize}
\textbf{Resultado (6 filas):} LJ\_003 y LJ\_005 lideran con 6 detecciones cada una; LJ\_004 solo
1. FRP m\'aximo m\'as alto: LJ\_005 (25.6 MW); FRP promedio m\'as alto: LJ\_003 (17.05 MW). \\
\textbf{Interpretaci\'on (anal\'itica):} LJ\_003 y LJ\_005 concentran m\'as detecciones que las
dem\'as celdas. LJ\_005 tiene el pico m\'aximo de FRP, pero LJ\_003 mantiene un promedio
ligeramente mayor, es decir, una intensidad m\'as consistente entre sus eventos. Es una
comparaci\'on entre celdas, no una tendencia temporal.

\subsection*{Consulta 7 -- GROUP BY + HAVING}
\textbf{Pregunta:} ¿qu\'e celdas y meses acumulan precipitaci\'on total por debajo de 20 mm?
\begin{scriptsize}
\begin{verbatim}
SELECT fc.celda_id, df.anio, df.mes, SUM(fc.precipitacion_diaria_mm) precip_mm,
       MIN(fc.precipitacion_diaria_mm) precip_min
FROM fact_clima fc JOIN dim_fecha df ON df.fecha_id = fc.fecha_id
GROUP BY fc.celda_id, df.anio, df.mes
HAVING SUM(fc.precipitacion_diaria_mm) < 20 ORDER BY precip_mm ASC;
\end{verbatim}
\end{scriptsize}
\textbf{Resultado representativo (9 filas en total; extremos):} LJ\_002/jul (0.00 mm, 1 d\'ia);
LJ\_001/jul (18.40 mm, 1 d\'ia, el m\'as cercano al umbral). Resultado completo en
\texttt{05\_resultados/resultados\_consultas.md}. \\
\textbf{Interpretaci\'on (anal\'itica):} 9 combinaciones celda-mes acumulan menos de 20 mm. La
mayor\'ia (7 de 9) tienen solo 1 o 2 d\'ias con registro en el mes, por lo que el acumulado bajo
puede reflejar tambi\'en cobertura de datos incompleta, no solo ausencia real de lluvia.

\subsection*{Consulta 8 -- CTE}
\textbf{Pregunta:} ¿qu\'e combinaciones celda-mes presentan a la vez detecciones y precipitaci\'on
mensual baja (< 20 mm)?
\begin{scriptsize}
\begin{verbatim}
WITH detecciones_mes AS (
    SELECT fi.celda_id, df.anio, df.mes, SUM(fi.numero_detecciones) det
    FROM fact_incendio fi JOIN dim_fecha df ON df.fecha_id = fi.fecha_id
    GROUP BY fi.celda_id, df.anio, df.mes),
clima_mes AS (
    SELECT fc.celda_id, df.anio, df.mes, SUM(fc.precipitacion_diaria_mm) precip
    FROM fact_clima fc JOIN dim_fecha df ON df.fecha_id = fc.fecha_id
    GROUP BY fc.celda_id, df.anio, df.mes)
SELECT d.celda_id, d.anio, d.mes, d.det, c.precip
FROM detecciones_mes d JOIN clima_mes c USING (celda_id, anio, mes)
WHERE d.det > 0 AND c.precip < 20 ORDER BY c.precip ASC;
\end{verbatim}
\end{scriptsize}
\textbf{Resultado (5 filas):} LJ\_003/jul (4 det., 2.10 mm); LJ\_005/ago (5 det., 4.50 mm);
LJ\_003/sep (2 det., 6.40 mm); LJ\_004/ago (1 det., 14.40 mm); LJ\_001/jul (2 det., 18.40 mm). \\
\textbf{Interpretaci\'on (anal\'itica):} en 5 combinaciones celda-mes coexisten detecciones y
precipitaci\'on mensual baja, destacando LJ\_003 en julio y LJ\_005 en agosto. Es una asociaci\'on
descriptiva observada en los datos de ejemplo, sin que constituya evidencia de causalidad ni deba
extrapolarse a otros periodos o celdas.

\subsection*{Consulta adicional -- fuera del m\'inimo exigido por la gu\'ia}
\textbf{Pregunta:} ¿qu\'e tipo de cobertura vegetal concentra m\'as detecciones? (usa la relaci\'on
N:M \texttt{dim\_celda} $\leftrightarrow$ \texttt{dim\_cobertura\_vegetal}).
\begin{scriptsize}
\begin{verbatim}
SELECT cv.nombre, COUNT(DISTINCT cc.celda_id) celdas, SUM(fi.numero_detecciones) detecciones
FROM dim_cobertura_vegetal cv
JOIN celda_cobertura cc ON cc.cobertura_id = cv.cobertura_id
JOIN fact_incendio fi ON fi.celda_id = cc.celda_id
GROUP BY cv.nombre ORDER BY detecciones DESC;
\end{verbatim}
\end{scriptsize}
\textbf{Resultado (5 filas):} Matorral (2 celdas, 10 det.); Pastizal (4 celdas, 10 det.); Bosque
nativo (2 celdas, 9 det.); Cultivo (2 celdas, 7 det.); \'Area urbana (1 celda, 2 det.). \\
\textbf{Interpretaci\'on (anal\'itica):} Matorral y Pastizal concentran el mayor n\'umero de
detecciones acumuladas, aunque Pastizal aparece en el doble de celdas que Matorral, por lo que su
detecci\'on promedio por celda es menor. Es una concentraci\'on descriptiva observada en los
datos de ejemplo; no implica que el tipo de cobertura sea causa de la ocurrencia de incendios.

% ==========================================================
% 6. VALIDACIONES
% ==========================================================
\section{Validaciones}

Ejecutadas contra \texttt{taller\_sql\_fireforest} (detalle completo en
\texttt{05\_resultados/\allowbreak validaciones.md} y en la secci\'on 5 de
\texttt{Grupo04\_Practica2\_SQL.sql}):

\begin{itemize}
    \item \textbf{Conteo de filas:} 82 en total (6+12+5+11+18+30 por tabla), como se dise\~n\'o.
    \item \textbf{PK \'unicas:} \texttt{fact\_incendio.incendio\_id} y \texttt{fact\_clima.clima\_id}
          sin duplicados (verificado comparando conteo total vs. conteo de valores distintos).
    \item \textbf{FK \'integras:} 0 filas de \texttt{fact\_incendio}/\texttt{fact\_clima} con
          celda o fecha inexistente.
    \item \textbf{Duplicados l\'ogicos:} 0 pares \texttt{(celda\_id, fecha\_id)} repetidos en
          \texttt{fact\_incendio}.
    \item \textbf{Valores negativos:} 0 precipitaciones y 0 valores de FRP negativos.
    \item \textbf{Cobertura $\leq$ 100\%:} 0 celdas superan el 100\% de cobertura acumulada por
          a\~no.
    \item \textbf{Coherencia de fechas:} de las 12 filas de \texttt{dim\_fecha}, 0 presentan
          \texttt{fecha\_id}, \texttt{anio} o \texttt{mes} incoherentes con \texttt{fecha}
          (verificado con consulta expl\'icita, no solo por la existencia del \texttt{CHECK}).
    \item \textbf{Pruebas de \texttt{CHECK}:} las 3 pruebas realizadas (mes fuera de rango,
          precipitaci\'on negativa, cobertura $>$100\%) fueron rechazadas correctamente por la
          base de datos.
\end{itemize}

% ==========================================================
% 7. CONCLUSIONES
% ==========================================================
\section{Conclusiones}

El modelo relacional dise\~nado permite integrar detecciones de incendio, precipitaci\'on y
cobertura vegetal por celda y fecha en una \'unica estructura consistente, con integridad
referencial verificada y sin necesidad de cruces manuales entre fuentes. La normalizaci\'on hasta
3FN (con la excepci\'on justificada de \texttt{dim\_fecha}) redujo la redundancia identificada en
el modelo inicial desnormalizado y facilit\'o el mantenimiento de los datos. Las 9 consultas
implementadas muestran que el esquema responde tanto preguntas operativas (filtrado puntual por
celda y fecha) como anal\'iticas (comparaciones entre celdas, meses y tipos de cobertura vegetal)
sin ambig\"uedad, siempre a partir de datos de ejemplo controlados.

Esta pr\'actica puede utilizarse como referencia o punto de partida conceptual para el componente
SQL del Proyecto Integrador FireForest, en tanto pr\'actica independiente: no modifica archivos,
esquema ni datos del Proyecto Integrador, y su incorporaci\'on eventual quedar\'ia sujeta a un
trabajo de integraci\'on posterior y expl\'icito.

\end{document}
